\subNsection{1}{Lessons Learned}{Alle}

\textbf{Kristian Popov}: Am meisten habe ich darüber gelernt, wie ich effizient
mit \gls{ki} arbeite.
Entscheidend war weniger die Wahl des Modells als seine Anbindung an das Projekt:
Über den Sprachserver (\gls{lsp}) und die Entwicklerwerkzeuge des Browsers fragte
das Modell Signaturen und Laufzeitverhalten ab, statt sie aus dem Quelltext
abzuleiten.
Zusammen mit einer Regeldatei für Konventionen und Architekturgrenzen brauchte es
weniger Tokens, und erfundene Bezeichner traten seltener auf.

Außerdem habe ich gelernt, Aufgaben konsequent zu planen und festzuhalten.
Über GitLab Issues und Milestones habe ich größere Vorhaben in überschaubare
Arbeitspakete zerlegt.
Gerade weil ich an zwei Themen parallel gearbeitet habe, blieb dadurch jederzeit
klar, was noch offen war.

Wie viel eine \gls{ci}-Pipeline abnimmt, habe ich ebenfalls erst im Projekt
gemerkt.
Formatierung, Lint, Tests und Build liefen bei jeder Änderung automatisch, sodass
ich Fehler nicht erst am Ende einer Entwicklung gefunden habe.
Gerade bei vielen kleinen Änderungen habe ich mich darauf verlassen können.

\textbf{Duy Tan Pham}: Im Projekt habe ich gelernt, größere Funktionen schrittweise zu entwickeln
und dabei ihre Einbindung in die gesamte Plattform mitzudenken. Durch einzelne
Ausbaustufen konnten neue Anforderungen aufgenommen und vorhandene Abläufe nach jeder
Erweiterung erneut geprüft werden. Der Abgleich mit den anderen Themenbereichen half
außerdem, die Bedienung einheitlich zu gestalten.

Bei Architekturentscheidungen zeigte sich, dass Wiederverwendung nicht automatisch mehr
Abstraktion bedeutet. Eine gemeinsame Zustandsabstraktion für mehrere Eingaben wurde
wieder entfernt, weil die eigenständigen Komponenten leichter nachzuvollziehen waren.
Gemeinsame Seitenrahmen und Bedienmuster waren dagegen an mehreren Stellen sinnvoll. Ich
habe daraus gelernt, eine Abstraktion nur dann einzuführen, wenn sie ein tatsächlich
wiederkehrendes Problem vereinfacht.

Eine weitere Erkenntnis war, dass nicht jede technisch mögliche Erweiterung das Produkt
verbessert. Einige Ansätze wirkten zunächst sinnvoll, passten nach genauerer Prüfung aber
nicht zu den Anforderungen oder erhöhten die Komplexität unnötig. Deshalb war es wichtig,
Entscheidungen erneut zu hinterfragen und Lösungen bei Bedarf bewusst zurückzunehmen.
Dadurch blieb der Fokus auf einem verständlichen und verlässlichen Lernerlebnis.

\textbf{Richard Schröder}: Für die Zusammenarbeit im Team war ein einheitliches Verständnis der Anforderungen entscheidend. Wenn vor der Umsetzung nicht genau geklärt war, wie eine Komponente aufgebaut sein soll, wichen die Teilergebnisse voneinander ab und mussten nachträglich mühsam angeglichen werden.

Beim Einsatz von \gls{ki}-Werkzeugen zeigte sich, dass der Zeitgewinn durch den hohen Aufwand für Prompt-Zerlegung, Fehlersuche sowie das Refactoring von totem oder doppeltem Code oft geringer ausfiel als erwartet. Da Modelle dazu neigten, Probleme unnötig zu verkomplizieren oder Änderungen an falschen Stellen im Projekt vorzunehmen, blieben gründliche manuelle Reviews und Bereinigungen notwendig. Zudem fiel der eigene Lerneffekt beim reinen Übernehmen generierter Lösungen geringer aus als beim eigenständigen Erarbeiten der Lösungswege.

Auch bei der Versionsverwaltung mit Git zeigte sich die Bedeutung enger Absprachen: Wurden Branches über längere Zeit ohne Zwischenabgleich entwickelt, führten gleichzeitige Änderungen an denselben Dateien schnell zu unübersichtlichen Merge-Konflikten. Kleinere und regelmäßig zusammengeführte Stände verhinderten diesen Mehraufwand.

\textbf{Enrico Ebert}: Bei der Arbeit mit \gls{bigint} wurde mir deutlich, dass JavaScripts mathematische Standardtypen bei großen Ganzzahlen an Präzision verlieren. Die Konvertierung zwischen verschiedenen Basen musste vollständig selbst implementiert werden, da keine native Funktion für Basis-zu-Basis-Umrechnung existiert. Die Visualisierung von \gls{ieee754} erwies sich als didaktisch besonders wertvoll: Studierende, die zuvor den Zusammenhang zwischen Dezimalzahl und Bitrepräsentation nicht verstanden, konnten durch die schrittweise Aufschlüsselung den Normalisierungsprozess nachvollziehen.

Die Validierung von \gls{utf8}-Byteströmen war komplexer als anfänglich angenommen. Insbesondere die korrekte Behandlung von Overlong Encodings erforderte intensives Lesen der \gls{unicode}-Spezifikation. Die Sicherheitsrelevanz dieser Prüfungen wurde mir erst durch die Auseinandersetzung mit Web-Sicherheitsthemen wie dem Facebook-Encoding-Bug bewusst. Die programmatische Übungsgenerierung erwies sich dagegen als besonders wertvoll: Sie bietet unbegrenzte Variationsmöglichkeiten ohne eine Datenbank pflegen zu müssen. Der Schwierigkeitsgrad wird über Parameter gesteuert und kann leicht angepasst werden.
